\section{Robustness}\label{sec:robustness}

Nine dimensions of robustness discipline the headline: strict invariance; cost-distribution estimator; kernel bandwidth; bid-level sample filter; temporal window; phased BEC enablement; alternative policy variants; bid-coordination diagnostics; and the Maskin--Riley FPSB upper bound on V3. Three claims survive each. \emph{(a)} The within-auction component remains the larger of the two in all six (class $\times$ window) cells; its share stays in the \withinShareLow--\withinShareHigh~percent range across cost-distribution estimators and across the two specifications of the post-policy SME pool. \emph{(b)} The welfare ranking V3 $\succ$ V0 in thick standardized non-pharmaceutical markets is preserved under cost-affiliation up to $\rho_c = \ipvRhoMax$, under the FPSB upper bound, and across both equilibrium-selection and strict-invariance treatments. \emph{(c)} The pharmaceutical bifurcation between the two specifications is itself the diagnostic of where modeling treatment of the protected pool is first-order. Detailed sensitivity grids, bandwidth and filter robustness, phased-adoption checks, alternative-variant comparisons, and bid-coordination screens are in the Online Appendix; this section summarizes the load-bearing diagnostics.

\subsection{Strict-primitive-invariance benchmark}

Table~\ref{tab:v3_strict_invariance} reports the limiting case in which $F_c^{\text{SME,Post}}$ is replaced by $F_c^{\text{SME,Pre}}$ while preserving the observed post-policy SME entry count---the strict-invariance bound on Assumption~\ref{a:setaside}. The simulated total price effect is $+\siDeltaTotalNp$ in non-pharma and $+\siDeltaTotalPh$ in pharma, with within-auction shares of \strictShareNp{} and \strictSharePh~percent. Welfare loss at $\lambda = \lambdaMain$ is \siWelfLossPctNp/\siWelfLossPctPh~percent of $p_{S_1}$.

What changes is the policy ranking in the thin heterogeneous pharma market. In non-pharma, the \optPrefThreshold{} price preference remains Pareto-superior under strict invariance as well as under the main specification. In pharma, the preference ranking reverses and the implied indifference weight (eq.~\ref{eq:welfare_weight}, Section~\ref{sec:discussion}) falls to $w^{\text{SME}}_\star = \welfWeightSiPh$ versus \welfWeightMainPh{} under the main specification. The bifurcation identifies precisely where equilibrium selection among protected bidders is first-order. Strict invariance imposes that protected-regime SMEs draw from the open-regime SME cost distribution---the most generous assumption to the policy, since it implies the protective regime selects entrants whose cost primitives match incumbents'. The realistic alternative is that protected entrants draw from the higher-cost tail (firms unprofitable under the open regime that become profitable under exclusion), which is what equilibrium selection in the main specification captures. Strict invariance therefore understates the welfare cost in pharma; the main-specification cost is the welfare-conservative number relative to the equilibrium-selection mechanism, not an aggressive upper bound.

\subsection{Cost-distribution regime, sample, and window}

Table~\ref{tab:v3_sensitivity_fc} runs the BNE decomposition under three estimators of $F_c^k$: all-bidders (baseline), losers-only (biases the upper tail upward), and Turnbull NPMLE (treats the winner as left-censored, delivers point identification). The three point estimates for $p_{S_3} - p_{S_1}$ are \costRegimeLosersNp/\costRegimeLosersPh{} (losers-only), \costRegimeAllNp/\costRegimeAllPh{} (all-bidders), and \costRegimeTurnNp/\costRegimeTurnPh{} (Turnbull) for non-pharma/pharma respectively. All cells lie within \turnbullDeviationPct~percent of the all-bidders estimate; the bootstrap CIs (Table~\ref{tab:v3_bootstrap_ci}) overlap materially across regimes. Turnbull raises the within-auction share in pharma from \bneShareIntensPh~percent to \turnPhShareUp~percent. A Gaussian-copula IPV relaxation at $\rho_c \in \ipvRhoGrid$ moves the within-auction share by less than \ipvShareDriftPp{} pp in either class; higher affiliation is inconsistent with the cross-modality consistency check (Table~\ref{tab:v3_cross_modality}).

The bid-level sample filter ($c_\epsilon \le \filterCepsUpper$, $n \ge \filterMinBidders$ baseline) varies along tight, very-tight, and loose cuts (Table~\ref{tab:v3_filter_sensitivity}); the baseline sits between tight and loose in both classes and is inside the bootstrap CI under any reasonable cut. The temporal window varies between 18-, 12-, and 6-month symmetric cuts (Table~\ref{tab:v3_window_sensitivity}). The non-pharmaceutical effect is essentially flat across windows; the pharmaceutical effect is monotonically increasing as the window shrinks, consistent with a transitional period in which the within-auction adjustment establishes more quickly than the participation-margin adjustment. The within-auction share varies by \windowShareDriftNp/\windowShareDriftPh{} pp across the three windows; the within-auction component remains larger in all six (class $\times$ window) cells. The 18m baseline is the mature, steady-state choice.

\subsection{Alternative policy variants and FPSB bound}\label{sec:fpsb_bound}

Table~\ref{tab:v3_apv} reports two intermediate variants alongside V0 and V3. \emph{V1} (partial 50/50 set-aside) yields $\Delta p$ at \vOneGapLo--\vOneGapHi~percent of V0's, roughly a third. \emph{V2} (full set-aside with the SME pool held at pre-period size) yields $\Delta p = $ \latentShockTimesNp/\latentShockTimesPh~percent of V0; without endogenous entry the participation margin absorbs \latentShockRatioNp/\latentShockRatioPh~percent of a heavier within-auction shock. V1--V3 span the policy-design space; V2 isolates entry's value, V1 prices a middle design, V3 captures the SME-protection value at near-zero price cost.

Section~\ref{sec:results} computes V3 under Vickrey equivalence; under asymmetric IPV, \citet{maskinriley2000} show that revenue equivalence between FPSB and ascending-clock does not generally hold. The FPSB--Vickrey gap on V3 (with $\rho = \prefRhoSmall$) is bounded above by $\rho \cdot \bar\sigma_c$ where $\bar\sigma_c$ is the SME-share-weighted CV of the UH-clean cost distributions ($\bar\sigma_c = \fpsbCovNp/\fpsbCovPh$; FPSB upper-bound gap $\fpsbBoundLoPct$--$\fpsbBoundHiPct~percent$ of $p^{\mathrm{ref}}$). Table~\ref{tab:maskin_riley_bound} reports the V0--V3 expected-price gap as \vZeroVThreeNp/\vZeroVThreePh{} in non-pharma/pharma---an order of magnitude larger than the FPSB upper-bound. The V3 $\succ$ V0 ranking is preserved at the FPSB upper bound in both classes; auction format is not a margin on which the comparison turns.

Phased-adoption checks (re-estimating the DiD with the BEC enablement window excluded), kernel-bandwidth sensitivity of the Convite GPV inversion, and four bid-coordination screens (Conley-style close-pair, Bajari--Ye persistent-pair, both conditioned on bidder count, bid-range quartiles, and CADMAT product class) all leave the within-auction and welfare-ranking conclusions intact. Per \citet{conley2016, kawainakabayashi2022} the bid-coordination evidence is suggestive rather than dispositive; the residual clustering after CADMAT conditioning is consistent with non-collusive mechanisms (subcontracting networks, joint ventures, regional logistics specialization) that the screens cannot separate from coordination. Full diagnostics in the Online Appendix.

\begin{table}[!htbp]
\centering
\begin{threeparttable}
\caption{Maskin and Riley (2000) bound on the FPSB-equivalent of the 10\% SME price preference (V3).}
\label{tab:maskin_riley_bound}
\small
\setlength{\tabcolsep}{4pt}
\begin{tabular}{lrrrrrrc}
\toprule
 & $p_{S_1}$ & $p_{V_0}$ & $p_{V_3}^{\text{Vick.}}$ & FPSB gap (UB) & $p_{V_3}^{\text{FPSB}}$ (UB) & $p_{V_0} - p_{V_3}^{\text{FPSB}}$ & V3 $\succ$ V0? \\
\midrule
Non-pharma & 0.759 & 1.018 & 0.755 & 0.045 & 0.800 & 0.218 & \textbf{Yes} \\
Pharma & 0.656 & 0.964 & 0.661 & 0.048 & 0.709 & 0.255 & \textbf{Yes} \\
\bottomrule
\end{tabular}
\begin{tablenotes}[flushleft]\footnotesize
\item Following \citet{maskinriley2000}, the asymmetric-IPV first-price-sealed-bid (FPSB) equivalent of the 10\% SME price preference V3 is not exactly Vickrey-equivalent; the expected-price gap is bounded above by a function of the type-specific coefficient-of-variation differential weighted by entry rates. Column ``FPSB gap (UB)'' reports this upper bound, computed as $\rho \cdot \bar\sigma_c^k$ where $\rho = 0.10$ is the preference parameter and $\bar\sigma_c^k$ is the SME-share-weighted average of the type-specific cost-distribution coefficient of variation in the stratum. Column ``$p_{V_3}^{\text{FPSB}}$ (UB)'' adds the bound to the Vickrey-equivalent $p_{V_3}$ reported in the preference-grid table (Online Appendix). Final column reports whether the V3 welfare-dominance over V0 is preserved at the FPSB upper bound; ``Yes'' indicates that even the worst-case FPSB price gap is not enough to reverse the policy ranking. This is a first-order bound; tighter bounds via the full ODE solver of \citet{lebrun1999} are reported as a future direction in Section~\ref{sec:discussion}.
\end{tablenotes}
\end{threeparttable}
\end{table}

